\documentclass[11pt]{article}
\usepackage{amsmath,amssymb,amsthm}
\usepackage[utf8]{inputenc}
\usepackage[T1]{fontenc}
\usepackage{amsmath,amssymb,amsthm}
\usepackage{fullpage}

\makeatletter
\renewcommand\section{\@startsection{section}{1}{\z@}
  {-3.5ex \@plus -1ex \@minus -.2ex}
  {2.3ex \@plus.2ex}
  {\normalfont\bfseries}}
\makeatother

\theoremstyle{plain}
\newtheorem{theorem}{Theorem}[section]
\newtheorem{lemma}[theorem]{Lemma}
\newtheorem{proposition}[theorem]{Proposition}
\newtheorem{corollary}[theorem]{Corollary}
\theoremstyle{definition}
\newtheorem{definition}[theorem]{Definition}
\newtheorem{example}[theorem]{Example}
\newtheorem{remark}[theorem]{Remark}
\newtheorem*{question}{Question}

\newcommand{\Z}{\mathbb{Z}}
\newcommand{\Hh}{\mathbb{H}}
\newcommand{\bp}{\boxplus}

\newcommand{\fptitle}{Review of Problem 5 -- ProofCouncil}
\newcommand{\reviewtag}{Reviewer 1}
\usepackage{fp-review}

\begin{document}
\maketitle

\section{Review}
{The solution is mathematically correct. It construct a counterexample using  infinite cyclic projective plane, a construction that already appears in the literature. The new observation is that by choosing parameters carefully, this give a negative answer to the question. This example is very similar to the example in other AI system submissions and is very different from the human solution. The proof writing here is too verbose at times: certain arguments can be omitted by pointing the reader to appropriate references.}

\paragraph{Correctness.}
\begin{itemize}
	\item  1: essentially flawless mathematical argument
	\item  0.5: minor errors that are easily corrected
	\item  0: blank solution, or acknowledged partial result
	\item  $-0.5$: major errors
	\item  $-1$: proposed proof makes no progress towards the result.
\end{itemize}
Note that points should be docked for wrong statements/ solutions.

\noindent Score: 1

\paragraph{Exposition.}
\begin{itemize}
	\item  1: well-written and clear solution.
	\item  0: argument that is possible to follow with serious effort
	\item  $-1$: very poor exposition
\end{itemize}

\noindent Score: 0

\paragraph{Attributions.}
\begin{itemize}
	\item  1: solutions contains fully appropriate attributions
		to related work and theorems being used.
	\item  0: moderate issues with missing attributions or hallucinated references
	\item  $-1$: blatant plagiarism or hallucinations
\end{itemize}

\noindent Score: 1


\newpage
\section{Submission with inline comments}


\begin{center}
{\textbf{A counterexample to the cubic implication for hyperfields}}
\end{center}

\medskip
\noindent\textbf{Abstract.}
In a field, $x+y+z=0$ implies $x^3+y^3+z^3-3xyz=0$.  We show that the
analogous implication for hyperfields,
\[
 0\in x\bp y\bp z
 \ \Longrightarrow\
 0\in x^3\bp y^3\bp z^3\bp(-xyz)\bp(-xyz)\bp(-xyz),
\]
fails in general, although it holds in every quotient $F/G$ of a field
by a multiplicative subgroup.  The counterexample is a hyperfield
attached, following Prenowitz and Connes--Consani, to an infinite
cyclic projective plane; the plane is built from a perfect difference
set in $\Z$ by an explicit variant of Hall's extension argument.  All
hyperfield axioms are verified directly, so the paper is
self-contained.

\section{Introduction}\label{sec:intro}

\paragraph{Conventions.}
A \emph{hyperoperation} on a set $\Hh$ is a map $\bp$ assigning to each
pair $(a,b)$ a nonempty subset $a\bp b\subseteq\Hh$.  It is extended to
subsets by the union convention
\begin{equation}\label{eq:union}
 A\bp B=\bigcup_{a\in A,\ b\in B}(a\bp b),
\end{equation}
an element being identified with its singleton.  Thus
$x\bp y\bp z$ means $(x\bp y)\bp z=\bigcup_{s\in x\bp y}(s\bp z)$, and
in an expression such as $(-xyz)\bp(-xyz)\bp(-xyz)$ each occurrence is a
separate summand; this is the reading intended in the question below.
In a hyperfield (Definition~\ref{def:hyperfield}) the hyperaddition is
associative, and then iterated sums do not depend on the bracketing
(Remark~\ref{rem:iterated}).

\paragraph{The question.}
In a field, $x+y+z=0$ implies $x^3+y^3+z^3-3xyz=0$, because of the
polynomial identity
\begin{equation}\label{eq:cubic}
 x^3+y^3+z^3-3xyz=(x+y+z)(x^2+y^2+z^2-xy-yz-zx).
\end{equation}
A \emph{hyperfield} is, roughly, a field whose addition is multivalued;
the precise definition is recalled in Section~\ref{sec:hyperfields}.
The question asks whether the implication survives in that generality.

\begin{question}
Let $\Hh$ be a hyperfield with hyperaddition $\bp$, and let
$x,y,z\in\Hh$ satisfy $0\in x\bp y\bp z$.  Does it follow that
\[
 0\in x^3\bp y^3\bp z^3\bp(-xyz)\bp(-xyz)\bp(-xyz)\,?
\]
\end{question}

\noindent
Here $x^3$ is the multiplicative cube and $-xyz$ the additive inverse
(the \emph{negative}) of the product $xyz$.  We call the implication
asked about the \emph{cubic implication}.  The answer is no.

\begin{theorem}\label{thm:main}
There exist a countably infinite commutative hyperfield $\Hh$ and
elements $x,y,z\in\Hh$ such that $0\in x\bp y\bp z$ but
\[
 0\notin x^3\bp y^3\bp z^3\bp(-xyz)\bp(-xyz)\bp(-xyz).
\]
\end{theorem}

\paragraph{Why a counterexample must be a non-quotient hyperfield.}
The most familiar hyperfields are Krasner's quotients $F/G$ of a field
$F$ by a subgroup $G$ of its multiplicative group
[Kr83] (recalled in Example~\ref{ex:quotient}).  In every such hyperfield the cubic
implication \emph{does} hold (Proposition~\ref{prop:quotient}): the
hypothesis $0\in x\bp y\bp z$ provides representatives $u,v,w\in F$
with $u+v+w=0$, and the field identity~\eqref{eq:cubic} for $u,v,w$
produces the required $0$.  Consequently a counterexample has to be a
hyperfield that is not a quotient of a field, so its hyperfield axioms
cannot be inherited from a field; we verify them directly in
Section~\ref{sec:hyperfield}, which is the main technical work of the
paper.

\paragraph{Strategy.}
The construction follows the geometric recipe of Prenowitz~[Pr43] and
Connes--Consani~[CC11, \S3].  Let $P$ be the point set of a projective
plane, in the weak sense that any two distinct points lie on exactly one
line and any two distinct lines meet in exactly one point, and write
$L(a,b)$ for the line through two distinct points $a,b$; three points
are \emph{collinear} if one line contains them and \emph{noncollinear}
otherwise, and the \emph{sides} of a noncollinear triple (a
\emph{triangle}) $a,b,c$ are the three lines $L(a,b),L(a,c),L(b,c)$.
On $\Hh=P\cup\{0\}$ define
\[
 \begin{aligned}
 a\bp b&=L(a,b)\setminus\{a,b\} &&(a,b\in P,\ a\neq b),\\
 a\bp a&=\{0,a\} &&(a\in P),\\
 0\bp a&=a\bp0=\{a\} &&(a\in\Hh)
 \end{aligned}
\]
(this is~\eqref{eq:addition} below).  If, in addition, every line is
infinite, $P$ is an abelian group, and
multiplication by every element of $P$ maps lines to lines, then $\Hh$
with this hyperaddition and the group multiplication is a commutative
hyperfield; this is proved for our plane in
Proposition~\ref{prop:hyperfield} by an argument that uses only these
properties (compare [CC11, Propositions~3.1 and~3.5], where four points
per line suffice).  Three features of this hyperfield drive the
counterexample.
\begin{enumerate}
\item[(a)] Every element is its own negative: $0\in a\bp b$ if and only
if $a=b$.  Hence $-xyz=xyz$, and for subsets $A,B\subseteq\Hh$ one has
$0\in A\bp B$ if and only if $A\cap B\neq\varnothing$
(Lemma~\ref{lem:zero}).
\item[(b)] If $a,b,c$ are distinct points on a common line $L$, then
$a\bp b\bp c=L\cup\{0\}$; in particular $0\in a\bp b\bp c$.
\item[(c)] If $a,b,c$ are noncollinear, then
$a\bp b\bp c=P\setminus\bigl(L(a,b)\cup L(a,c)\cup L(b,c)\bigr)$ is
the complement of the three sides of the triangle $a,b,c$
(Lemma~\ref{lem:triangle}); in particular $0\notin a\bp b\bp c$.
\end{enumerate}
Now let $x,y,z$ be distinct collinear points, so that $0\in x\bp y\bp z$
by~(b), and put $q=xyz=-xyz$.  The defining rules give
$q\bp q\bp q=\{0,q\}\bp q=\{q\}\cup\{0,q\}=\{0,q\}$, so by associativity
the six-term sum of the question equals $S\bp\{0,q\}$ with
$S=x^3\bp y^3\bp z^3$, and by~(a) it contains $0$ if and only if $S$
meets $\{0,q\}$.  Suppose the cubes $x^3,y^3,z^3$ are noncollinear.
Then (c) says that $0\notin S$ and that $q\in S$ exactly when $q$ lies
on none of the sides of the triangle $x^3,y^3,z^3$.  The question has
thus become one of pure incidence geometry: the cubic implication fails
for $x,y,z$ precisely when the product $q=xyz$ lies on a side of the
``cube triangle''.  So we need a plane whose points form a group acting
by collineations, containing three collinear points whose cubes are
noncollinear and whose product lies on a side of the cube triangle.

\paragraph{Cyclic planes.}
Let $P=\{t^j:j\in\Z\}$ be an infinite cyclic group and let
$D\subseteq\Z$ be a \emph{perfect difference set}: every nonzero
integer has exactly one representation $d-e$ with $d,e\in D$.  The sets
$L_k=\{t^{k+d}:d\in D\}$, $k\in\Z$, are then the lines of a projective
plane on $P$, and multiplication by $t^r$ carries $L_k$ to $L_{k+r}$
(Section~\ref{sec:plane}).  This is the classical correspondence between
cyclic planes and difference sets, due to Singer~[Si38] for finite
planes and to Hall~[Ha47] for the group $\Z$.  Choose
\[
 x=1,\qquad y=t,\qquad z=t^3,\qquad\text{so that}\qquad
 x^3=1,\quad y^3=t^3,\quad z^3=t^9,\quad q=xyz=t^4.
\]
The requirements become conditions on $D$: if $0,1,3\in D$ then $x,y,z$
lie on $L_0$; if $20,24,29\in D$ then $1$, $q=t^4$ and $t^9$ lie on
$L_{-20}$, which is therefore the side $L(1,t^9)$ of the cube triangle;
and the uniqueness of differences then forces $23\notin D$ (because
$23-20=3-0$), so that $t^3\notin L_{-20}$ and the cubes are
noncollinear.  Hall proved that every finite set of integers with
pairwise distinct differences is contained in a perfect difference set
of $\Z$ [Ha47, Theorem~3.1].  Section~\ref{sec:diffset} carries this
out explicitly, with a concrete spacing rule, for the seed
\[
 A_0=\{0,1,3,20,24,29\},
\]
whose fifteen positive differences are distinct.  (The value $20$ is one
convenient choice making the differences of the seed distinct; nothing
else about it is used.)

\paragraph{Organization.}
Section~\ref{sec:hyperfields} recalls hyperfields and quotient
hyperfields and proves that the cubic implication holds in the latter.
Section~\ref{sec:diffset} constructs the perfect difference set $D$,
Section~\ref{sec:plane} the cyclic plane, and
Section~\ref{sec:hyperfield} the hyperfield $\Hh$, including the
verification of all axioms.  Section~\ref{sec:counterexample} proves
Theorem~\ref{thm:main} and deduces that $\Hh$ is not a quotient of a
field.  Credits for the ingredients are collected in
Section~\ref{sec:credits}.

\section{Hyperfields and quotient hyperfields}\label{sec:hyperfields}

\begin{definition}[Krasner {[Kr83]}]\label{def:hyperfield}
A \emph{commutative hyperfield} is a set $\Hh$ with two distinct
distinguished elements $0$ and $1$, a multiplication making
$\Hh\setminus\{0\}$ an abelian group with identity $1$ and satisfying
$0\cdot a=a\cdot 0=0$, and a hyperoperation $\bp$ (the
\emph{hyperaddition}) such that for all $a,b,c\in\Hh$:
\begin{enumerate}
\item[(H1)] $a\bp b=b\bp a$;
\item[(H2)] $(a\bp b)\bp c=a\bp(b\bp c)$, both sides being read with the
convention~\eqref{eq:union};
\item[(H3)] $0\bp a=\{a\}$;
\item[(H4)] there is a unique element $-a\in\Hh$ with $0\in a\bp(-a)$;
\item[(H5)] (reversibility) if $c\in a\bp b$, then $b\in c\bp(-a)$;
\item[(H6)] $a(b\bp c)=ab\bp ac$, where $a(b\bp c)=\{ad:d\in b\bp c\}$.
\end{enumerate}
\end{definition}

A set with a hyperoperation and an element $0$ satisfying (H1)--(H5) is
called a \emph{canonical hypergroup}, a term due to Mittas~[Mi72]; so a
hyperfield is a canonical hypergroup with a compatible multiplication.
A field is a hyperfield whose sums are singletons.  Hyperfields with noncommutative
multiplication are also considered in the literature; since the
question concerns an arbitrary hyperfield, a commutative counterexample
suffices, and ``hyperfield'' will mean commutative hyperfield below.

\begin{remark}[iterated sums]\label{rem:iterated}
The convention~\eqref{eq:union} makes $\bp$ a binary operation on the
nonempty subsets of $\Hh$, and (H1), (H2) make this operation
commutative and associative: for nonempty $A,B,C$,
\[
 (A\bp B)\bp C
 =\bigcup_{a\in A,\,b\in B,\,c\in C}(a\bp b)\bp c
 =\bigcup_{a\in A,\,b\in B,\,c\in C}a\bp(b\bp c)
 =A\bp(B\bp C).
\]
By the usual generalized associativity for an associative commutative
binary operation, an iterated sum $x_1\bp\cdots\bp x_n$ is independent
of the bracketing and of the order of the summands, and may be regrouped
freely.  In Section~\ref{sec:counterexample} we use this to write a
six-term sum as $(x^3\bp y^3\bp z^3)\bp(q\bp q\bp q)$.  Note also
$A\bp\{0\}=A$, by (H3).
\end{remark}

\begin{lemma}\label{lem:zero}
Let $\Hh$ be a hyperfield and $A,B\subseteq\Hh$ nonempty.  Then
$0\in A\bp B$ if and only if there is $a\in A$ with $-a\in B$.
In particular, if every element of $\Hh$ satisfies $-a=a$, then
$0\in A\bp B$ if and only if $A\cap B\neq\varnothing$.
\end{lemma}
\begin{proof}
For elements, $0\in a\bp b$ holds if and only if $b=-a$, by the
uniqueness in (H4).  Now take the union over $a\in A$, $b\in B$.
\end{proof}

The second statement is the criterion of [CC11, Lemma~3.2], stated
there for vector spaces over the Krasner hyperfield $\mathbf K$, in
which likewise every element is its own negative.

\begin{example}[quotient hyperfields, Krasner {[Kr83]}]\label{ex:quotient}
Let $F$ be a field and $G\leq F^\times$ a subgroup of its multiplicative
group.  Let $F/G=\{0\}\cup\{uG:u\in F^\times\}$ be the set of orbits of
$G$ acting on $F$ by multiplication, and let $\pi\colon F\to F/G$ be the
orbit map, $\pi(0)=0$ and $\pi(u)=uG$ for $u\neq0$.  Multiplication on
$F/G$ is defined by $\pi(u)\pi(v)=\pi(uv)$, which is well defined, and
hyperaddition by
\[
 \pi(u)\bp\pi(v)
 =\bigl\{\pi(u'+v'):\ u'\in\pi^{-1}(\pi(u)),\ v'\in\pi^{-1}(\pi(v))\bigr\},
\]
the set of orbits of all sums of a representative of the first class
and a representative of the second.  Krasner~[Kr83] proved that $F/G$
is a commutative hyperfield.  Its negatives are $-\pi(u)=\pi(-u)$,
because $0=\pi(u+(-u))\in\pi(u)\bp\pi(-u)$ and negatives are unique by
(H4).  The fibres $\pi^{-1}(h)$ are $\{0\}$ and the cosets $uG$; they
are stable under multiplication by elements of $G$.
\end{example}

\begin{proposition}\label{prop:quotient}
Let $F$ be a field, $G\leq F^\times$ a subgroup, and $x,y,z\in F/G$.
If $0\in x\bp y\bp z$, then
$0\in x^3\bp y^3\bp z^3\bp(-xyz)\bp(-xyz)\bp(-xyz)$.
\end{proposition}
\review{This is not a part of the main argument, it might be better to move this to the end instead.}
\begin{proof}
We use two general facts about the orbit map $\pi$.

\smallskip\noindent
(i) \emph{If $u_1,\dots,u_n\in F$, then
$\pi(u_1+\cdots+u_n)\in\pi(u_1)\bp\cdots\bp\pi(u_n)$.}
For $n=1$ there is nothing to prove.  For $n\geq2$, the element
$u_1+\cdots+u_{n-1}$ is a representative of its own class and $u_n$ of
its own, so by the definition of $\bp$,
\[
 \pi(u_1+\cdots+u_n)\in\pi(u_1+\cdots+u_{n-1})\bp\pi(u_n),
\]
and $\pi(u_1+\cdots+u_{n-1})\in\pi(u_1)\bp\cdots\bp\pi(u_{n-1})$ by
induction; the convention~\eqref{eq:union} combines the two memberships.

\smallskip\noindent
(ii) \emph{If $0\in x\bp y\bp z$, then there are representatives
$u\in\pi^{-1}(x)$, $v\in\pi^{-1}(y)$, $w\in\pi^{-1}(z)$ with
$u+v+w=0$.}  By~\eqref{eq:union}, $0\in s\bp z$ for some $s\in x\bp y$.
By the definition of $\bp$, $0=\pi(s'+w)$ for some $s'\in\pi^{-1}(s)$
and $w\in\pi^{-1}(z)$; as $\pi^{-1}(0)=\{0\}$, this means $s'=-w$.
Also $s=\pi(u'+v')$ for some $u'\in\pi^{-1}(x)$, $v'\in\pi^{-1}(y)$.
If $u'+v'\neq0$, then $s'$ and $u'+v'$ lie in the same coset of $G$, so
$s'=(u'+v')g$ for some $g\in G$; put $u=u'g$ and $v=v'g$.  These still
represent $x$ and $y$, because fibres are $G$-stable, and
$u+v+w=s'+w=0$.  If $u'+v'=0$, then $s=0$, hence $s'=0$ and $w=0$, and
$u=u'$, $v=v'$ work.

\smallskip
Now let $0\in x\bp y\bp z$ and choose $u,v,w$ as in (ii).  The
identity~\eqref{eq:cubic} in the field $F$ gives
\[
 u^3+v^3+w^3+(-uvw)+(-uvw)+(-uvw)=0.
\]
Since $\pi$ is multiplicative and $\pi(-a)=-\pi(a)$
(Example~\ref{ex:quotient}), the six summands represent
$\pi(u^3)=x^3$, $\pi(v^3)=y^3$, $\pi(w^3)=z^3$ and
$\pi(-uvw)=-xyz$ (three times).  Applying (i) to them gives
\[
 0=\pi(0)\in x^3\bp y^3\bp z^3\bp(-xyz)\bp(-xyz)\bp(-xyz).\qedhere
\]
\end{proof}

An \emph{isomorphism} of hyperfields is a bijection $f$ preserving $0$,
multiplication and hyperaddition (that is, $f(a\bp b)=f(a)\bp f(b)$ as
sets); it transports the cubic implication verbatim.  Together with
Theorem~\ref{thm:main}, Proposition~\ref{prop:quotient} will therefore
show that the hyperfield constructed below is not isomorphic to any
quotient $F/G$ (Corollary~\ref{cor:nonquotient}).

\section{A perfect difference set in the integers}\label{sec:diffset}

\begin{definition}\label{def:pds}
A subset $D\subseteq\Z$ is a \emph{perfect difference set} if the map
\[
 \{(d,e)\in D\times D:\ d\neq e\}\longrightarrow\Z\setminus\{0\},
 \qquad (d,e)\longmapsto d-e,
\]
is bijective; that is, every nonzero integer has exactly one expression
$d-e$ with $d,e\in D$.  For a finite set $A\subseteq\Z$ write
\[
 \Delta(A)=\{b-a:\ a,b\in A,\ a<b\}
\]
for its set of positive differences, and say that $A$ has
\emph{distinct differences} if the expressions $b-a$, for $a<b$ in $A$,
are pairwise distinct, i.e.\ if $|\Delta(A)|=\binom{|A|}{2}$.
\end{definition}

The terminology is Hall's~[Ha47]; the finite analogue, difference sets
modulo $n^2+n+1$, goes back to Singer~[Si38].  Recall from the
introduction that, by [Ha47, Theorem~3.1], every finite set with
distinct differences is contained in a perfect difference set of $\Z$.
We now carry this out for a specific seed and with an explicit spacing
rule; the result is Lemma~\ref{lem:diff}.
\review{It suffices to cite [Ha47, Theorem~3.1] instead of reproducing Hall's argument.}

\paragraph{The seed.}
Let $A_0=\{0,1,3,20,24,29\}$.  Its positive differences, listed by
smaller element, are
\[
 \begin{array}{ll}
 \text{from }0:& 1,\ 3,\ 20,\ 24,\ 29,\\
 \text{from }1:& 2,\ 19,\ 23,\ 28,\\
 \text{from }3:& 17,\ 21,\ 26,\\
 \text{from }20:& 4,\ 9,\\
 \text{from }24:& 5,
 \end{array}
\]
so that, in increasing order,
\[
 \Delta(A_0)=\{1,2,3,4,5,9,17,19,20,21,23,24,26,28,29\}.
\]
\review{This part is too verbose.}
These are $15=\binom62$ distinct numbers, so $A_0$ has distinct
differences.  The roles of the seed elements were explained in the
introduction: $0,1,3$ make $x,y,z$ collinear, and $20,24,29$ make
$1,t^4,t^9$ collinear.

\paragraph{The extension step.}
Suppose a finite set $A_n\subseteq\Z_{\geq0}$ with distinct differences
has been constructed.  Define
\begin{equation}\label{eq:parameters}
 M_n=\max A_n,\qquad
 m_n=\min\bigl(\Z_{>0}\setminus\Delta(A_n)\bigr),\qquad
 N_n=2M_n+2m_n+1,
\end{equation}
and set
\begin{equation}\label{eq:recursion}
 A_{n+1}=A_n\cup\{N_n,\ N_n+m_n\},\qquad
 D=\bigcup_{n\geq0}A_n.
\end{equation}
The integer $m_n$, the least positive difference not yet realized,
exists because $A_n$ is finite.  The idea of the step is simple: the two
new elements realize the missing difference $m_n$, and they are placed
so far to the right of $A_n$ that every other new difference exceeds
every old one.  For instance, $m_0=6$ and $N_0=2\cdot29+2\cdot6+1=71$,
so $A_1=A_0\cup\{71,77\}$; then $m_1=7$ and $N_1=169$, so
$A_2=A_1\cup\{169,176\}$.

\begin{lemma}\label{lem:diff}
Every $A_n$ has distinct differences.  The set $D$ is an infinite
perfect difference set with $D\subseteq\Z_{\geq0}$ and $\min D=0$.
\end{lemma}
\begin{proof}
\emph{Step 1: $A_{n+1}$ has distinct differences.}
Abbreviate $A_n,M_n,m_n,N_n$ by $A,M,m,N$.  All elements of $A$ lie in
$[0,M]$, so every old difference is at most $M$.  The positive
differences of $A_{n+1}=A\cup\{N,N+m\}$ are the old ones, the
difference $(N+m)-N=m$ of the two new elements, and the mixed
differences
\[
 N-a\quad\text{and}\quad N+m-a\qquad(a\in A),
\]
which are positive because both new elements exceed $M$.  Since
$a\leq M$,
\[
 N-a\geq N-M=M+2m+1>\max\{M,m\},
\]
and the same bound holds for $N+m-a$.  Consequently no mixed difference
coincides with an old difference or with $m$, and $m$ is not an old
difference by its choice.  Within the family $\{N-a\}$, equal values
force equal $a$, and likewise within $\{N+m-a\}$.  A coincidence
$N-a=N+m-b$ between the two families would give $b-a=m$ with
$a,b\in A$, i.e.\ $m\in\Delta(A)$, contradicting the choice of $m$.
Thus all positive differences of $A_{n+1}$ are distinct, and by
induction every $A_n$ has distinct differences.

\emph{Step 2: every nonzero integer is a difference.}
By the minimality of $m_n$, the integers $1,\dots,m_n-1$ lie in
$\Delta(A_n)\subseteq\Delta(A_{n+1})$, and $m_n\in\Delta(A_{n+1})$ by
construction.  Hence $\{1,\dots,m_n\}\subseteq\Delta(A_{n+1})$ and
$m_{n+1}>m_n$.  The sequence $(m_n)$ is therefore unbounded, so every
positive integer $k$ satisfies $k<m_n$ for some $n$ and then
$k\in\Delta(A_n)$; thus $k=d-e$ with $d,e\in D$.  For negative $k$,
reverse the pair.

\emph{Step 3: uniqueness.}
Suppose $k=d-e=d'-e'$ with $(d,e)\neq(d',e')$ in $D\times D$ and, say,
$k>0$.  The sets $A_n$ increase with $n$, so $d,e,d',e'$ all lie in one
$A_n$, and $k$ would have two different expressions as a positive
difference in $A_n$, contrary to Step~1.  For $k<0$ negate both
expressions.

\emph{Step 4.}  Each step adjoins two elements larger than $M_n$, so $D$
is infinite; and $0\in A_0\subseteq D\subseteq\Z_{\geq0}$.
\end{proof}

\section{The cyclic plane}\label{sec:plane}

Let
\[
 P=\{t^j:\ j\in\Z\}
\]
be an infinite cyclic group with generator $t$, written
multiplicatively: $t^it^j=t^{i+j}$, with identity $1=t^0$, and distinct
exponents denote distinct elements.  Using the perfect difference set
$D$ of Lemma~\ref{lem:diff}, define for every $k\in\Z$ the \emph{line}
\[
 L_k=\{t^{k+d}:\ d\in D\}=t^k L_0 .
\]
Thus the lines are the translates of the single line $L_0=\{t^d:d\in D\}$.
Three points are \emph{collinear} if some line contains all of them, and
\emph{noncollinear} otherwise; a noncollinear triple $a,b,c$ is also
called a \emph{triangle} with \emph{vertices} $a,b,c$, and once (I1)
below is available its \emph{sides} are the lines
$L(a,b),L(a,c),L(b,c)$.  The following lemma collects everything that
will be used about the plane.

\begin{lemma}\label{lem:incidence}
The points of $P$ and the lines $L_k$ have the following properties.
\begin{enumerate}
\item[(I1)] Any two distinct points of $P$ lie on exactly one line.
\item[(I2)] Any two distinct lines meet in exactly one point.
\item[(I3)] The lines $L_k$, $k\in\Z$, are pairwise distinct.
\item[(I4)] Every line is infinite.  (Also, no line is all of $P$, since
$t^{k-1}\notin L_k$; this will not be needed.)
\item[(I5)] $t^rL_k=L_{k+r}$ for all $r,k\in\Z$.
\end{enumerate}
\end{lemma}
\review{Attributions are missing for this lemma. They appear in the introduction and in Section 9, but should be here as well.}
\begin{proof}
(I3) Since $\min D=0$, the least exponent occurring in $L_k$ is $k$.

(I1) Let $t^i\neq t^j$, so $i\neq j$.  The point $t^i$ lies on $L_k$ if
and only if $i-k\in D$.  Hence $L_k$ contains both points if and only if
$i-k=d$ and $j-k=e$ for some $d,e\in D$, and then $d-e=i-j\neq0$.  By
Lemma~\ref{lem:diff} there is exactly one pair $(d,e)\in D\times D$ with
$d-e=i-j$, so exactly one $k$, namely $k=i-d=j-e$, works; by (I3) this
means that exactly one line contains both points.

(I2) Let $k\neq\ell$.  A point $t^j$ lies on both $L_k$ and $L_\ell$ if
and only if $j-k=d\in D$ and $j-\ell=e\in D$, and then
$d-e=\ell-k\neq0$.  Again there is exactly one such pair $(d,e)$, hence
exactly one such $j=k+d$.

(I4) $D$ is infinite, and $-1\notin D$.

(I5) $t^r L_k=\{t^{r+k+d}:d\in D\}=L_{k+r}$.
\end{proof}

By (I1) we may write $L(a,b)$ for the unique line through two distinct
points $a,b$; thus, for $c\neq a,b$, the points $a,b,c$ are collinear if
and only if $c\in L(a,b)$.  A useful consequence of (I5): for distinct
$a,b\in P$ and $u\in P$,
\begin{equation}\label{eq:translation}
 uL(a,b)=L(ua,ub),
\end{equation}
because $uL(a,b)$ is a line by (I5) and contains the distinct points
$ua,ub$.  Properties (I1) and (I2) say that $P$ with the lines $L_k$ is
a projective plane in the weak sense used in the introduction, and (I5)
says that the group $P$ acts on it by collineations.  Projective planes admitting a
cyclic group of collineations that acts sharply transitively on the
points (for any two points exactly one group element carries the first
to the second) are the \emph{cyclic planes} of Singer and Hall; in the
terminology of [CC11, \S3], a group structure on the point set for
which multiplications are collineations, such as our $P$, is a
two-sided incidence group of the plane.

\section{The hyperfield}\label{sec:hyperfield}

Adjoin to $P$ a new element $0\notin P$ and put $\Hh=P\cup\{0\}$.
Extend the multiplication of $P$ by $0a=a0=0$ for all $a\in\Hh$; the
multiplicative identity is $1=t^0$, which is different from the new
element $0$.  Define the hyperaddition by
\begin{equation}\label{eq:addition}
 \begin{aligned}
 0\bp a&=a\bp0=\{a\} &&(a\in\Hh),\\
 a\bp a&=\{0,a\} &&(a\in P),\\
 a\bp b&=L(a,b)\setminus\{a,b\}
       &&(a,b\in P,\ a\neq b).
 \end{aligned}
\end{equation}
In words: the sum of two distinct points consists of the other points of
the line joining them, and the sum of a point with itself is the point
together with $0$.  This is Prenowitz's hyperaddition on a projective
geometry [Pr43], in the form used by Connes and Consani
[CC11, Proposition~3.1], who show that it yields a canonical hypergroup
whenever every line has at least four points, and that a two-sided
incidence group turns it into a hyperfield [CC11, Proposition~3.5].  We
do not rely on those results but verify the axioms directly; this also
makes visible exactly which incidence properties are used, namely
(I1), (I2), (I4) and~\eqref{eq:translation}.

Throughout this section we write, for $a,b,c\in\Hh$,
\[
 T(a,b,c)=(a\bp b)\bp c=\bigcup_{d\in a\bp b}(d\bp c)
\]
for the left-bracketed triple sum.  The first lemma is the heart of the
associativity proof: for a triangle the triple sum has a description
that is visibly symmetric in the three summands.  The formula also
appears in the proof of [CC11, Proposition~3.1].

\begin{lemma}[triangle formula]\label{lem:triangle}
For three noncollinear points $a,b,c\in P$,
\begin{equation}\label{eq:triangle}
 T(a,b,c)=(a\bp b)\bp c
 =P\setminus\bigl(L(a,b)\cup L(a,c)\cup L(b,c)\bigr).
\end{equation}
In particular $0\notin T(a,b,c)$, and $T(a,b,c)$ is disjoint from the
three sides of the triangle $abc$.
\end{lemma}
\begin{proof}
Noncollinear points are distinct, and $c\notin L(a,b)$,
$b\notin L(a,c)$, $a\notin L(b,c)$; so the three sides are three distinct
lines, and by (I2) any two of them meet only in their common vertex.

\emph{The left side is contained in the right side.}
Let $r\in T(a,b,c)$.  Then $r\in d\bp c$ for some
$d\in a\bp b=L(a,b)\setminus\{a,b\}$.  Since $c\notin L(a,b)$ we have
$d\neq c$, so $d\bp c=L(d,c)\setminus\{d,c\}$; in particular $r\neq0$,
$r\in P$, and $r\neq d,c$.  The line $L(d,c)$ is different from
$L(a,b)$ (it contains $c$), so it meets $L(a,b)$ only in $d$, by (I2);
as $r\neq d$, $r\notin L(a,b)$.  Next, $d\notin L(a,c)$: the lines
$L(a,c)$ and $L(a,b)$ are distinct and meet only in $a$, while
$d\in L(a,b)$ and $d\neq a$.  Hence $L(d,c)\neq L(a,c)$, these two
lines meet only in $c$, and $r\neq c$ gives $r\notin L(a,c)$.  The same
argument with $b$ in place of $a$ gives $r\notin L(b,c)$.

\emph{The right side is contained in the left side.}
Let $r\in P$ lie on none of the three sides.  Then $r\neq c$, and the
line $L(c,r)$ is different from $L(a,b)$ because $c\notin L(a,b)$; by
(I2) the two lines meet in exactly one point $d$.  If $d$ were equal to
$a$, then $a\in L(c,r)$ would force $L(c,r)=L(a,c)$ by (I1), placing
$r$ on the side $L(a,c)$; so $d\neq a$, and likewise $d\neq b$.  Also
$d\neq c$ (as $c\notin L(a,b)$) and $d\neq r$ (as $r\notin L(a,b)$).
Therefore $d\in L(a,b)\setminus\{a,b\}=a\bp b$, and
$r\in L(c,r)\setminus\{d,c\}=L(d,c)\setminus\{d,c\}=d\bp c$, so
$r\in T(a,b,c)$.
\end{proof}

\begin{proposition}\label{prop:hyperfield}
With the multiplication and hyperaddition above, $\Hh$ is a commutative
hyperfield in the sense of Definition~\ref{def:hyperfield}.  Moreover
\begin{equation}\label{eq:zero}
 0\in a\bp b\quad\Longleftrightarrow\quad a=b\qquad(a,b\in\Hh),
\end{equation}
so every element is its own negative, $-a=a$.
\end{proposition}
\begin{proof}
All sums are nonempty: a line minus two points is nonempty by (I4).
Commutativity (H1) and the neutrality of $0$ (H3) are built into
\eqref{eq:addition}.

\emph{Negatives (H4).}  Inspecting the three rules of
\eqref{eq:addition} gives~\eqref{eq:zero}: $0$ lies in $a\bp a$ for
every $a$ (including $a=0$), in $0\bp b=\{b\}$ only if $b=0$, and in
$L(a,b)\setminus\{a,b\}\subseteq P$ never.  Thus $a$ itself is the
unique element $-a$ with $0\in a\bp(-a)$.

\emph{Reversibility (H5).}  Since $-a=a$, we must show that
$c\in a\bp b$ implies $b\in c\bp a$; we show in fact that the two
conditions
\[
 c\in a\bp b\qquad\text{and}\qquad b\in c\bp a
\]
are equivalent.  If one of $a,b,c$ is $0$, both conditions say that the
other two are equal: for $a=0$ by (H3), for $b=0$ by (H3)
and~\eqref{eq:zero}, and for $c=0$ by~\eqref{eq:zero} and (H3).  Let
$a,b,c\in P$.  If $a=b=c$, both conditions hold ($a\in a\bp a$).  If
exactly two of them are equal, both conditions fail: a sum
$u\bp u=\{0,u\}$ never contains a point different from $u$, and a sum
$u\bp v$ with $u\neq v$ contains neither $u$ nor $v$.  If $a,b,c$ are
distinct, both conditions say that $a,b,c$ are collinear, by (I1).

\emph{Associativity (H2).}  The commutativity part of
Remark~\ref{rem:iterated} uses only (H1), so
\[
 a\bp(b\bp c)=(b\bp c)\bp a=T(b,c,a).
\]
Hence (H2) follows once we show that $T(a,b,c)$ depends only on the
multiset $\{a,b,c\}$ and not on the order of the summands.  A triple
either contains $0$, or lies in $P^3$ and has three equal summands,
exactly two equal summands, or three distinct summands, the last being
collinear or not; we compute $T$ in each of these five cases.
\begin{itemize}
\item \emph{A zero summand.}  By (H3) and $A\bp\{0\}=A$,
$T(0,a,b)=T(a,0,b)=T(a,b,0)=a\bp b$, which is symmetric in $a,b$.
\item \emph{Three equal nonzero summands $a,a,a$.}
$T(a,a,a)=\{0,a\}\bp a=(0\bp a)\cup(a\bp a)=\{0,a\}$.
\item \emph{Exactly two equal summands.}  Let the summands be $a,a,b$
with $a\neq b$ in $P$, in any order, and put $L=L(a,b)$.  Then
\[
 T(a,a,b)=\{0,a\}\bp b=\{b\}\cup\bigl(L\setminus\{a,b\}\bigr)
 =L\setminus\{a\},
\]
and, since $d\bp a=L\setminus\{d,a\}$ for $d\in L\setminus\{a,b\}$,
\[
 T(a,b,a)=T(b,a,a)
 =\bigcup_{d\in L\setminus\{a,b\}}\bigl(L\setminus\{a,d\}\bigr)
 =L\setminus\{a\}.
\]
For the last equality: the union is contained in $L\setminus\{a\}$, and
a point $r\in L\setminus\{a\}$ is retained by choosing
$d\in L\setminus\{a,b,r\}$, which is possible because $L$ is infinite
(I4).  So all three orderings give $L(a,b)\setminus\{a\}$.
\item \emph{Three distinct collinear summands} $a,b,c$ on a line $L$.
Here $a\bp b=L\setminus\{a,b\}$, and for $d\in L\setminus\{a,b\}$ the
summand $d\bp c$ is $\{0,c\}$ if $d=c$ and $L\setminus\{d,c\}$ if
$d\neq c$ (because then $L(d,c)=L$).  Hence
\[
 T(a,b,c)=\{0,c\}\cup\bigcup_{d\in L\setminus\{a,b,c\}}
 \bigl(L\setminus\{d,c\}\bigr)=\{0,c\}\cup\bigl(L\setminus\{c\}\bigr)
 =L\cup\{0\},
\]
where a point $r\in L\setminus\{c\}$ is retained by choosing
$d\in L\setminus\{a,b,c,r\}$, again possible by (I4).  The result
$L\cup\{0\}$ depends only on the line $L$, hence only on the multiset
$\{a,b,c\}$.
\item \emph{Three noncollinear summands.}  By
Lemma~\ref{lem:triangle}, $T(a,b,c)$ is the complement in $P$ of the
union of the three sides, which is symmetric in $a,b,c$.
\end{itemize}
This proves (H2).  Thus $(\Hh,\bp,0)$ is a canonical hypergroup.  (It
also satisfies $a\bp\Hh=\Hh$, a condition some authors include: given
$b$, pick $c\in b\bp a$; then $b\in c\bp a$ by (H5).)

\emph{Multiplication and distributivity (H6).}  The nonzero elements
form the abelian group $P$ with identity $1\neq0$, and $0$ is
absorbing.  Let $u\in P$.  For distinct $a,b\in P$,
\eqref{eq:translation} gives
\[
 u(a\bp b)=uL(a,b)\setminus\{ua,ub\}=L(ua,ub)\setminus\{ua,ub\}
 =ua\bp ub,
\]
using that multiplication by $u$ is a bijection of $P$.  For equal
nonzero summands, $u\{0,a\}=\{0,ua\}=ua\bp ua$; if a summand is $0$,
both sides of (H6) are the singleton of the product of $u$ with the
other summand.  Finally, for $u=0$ both sides of (H6) equal $\{0\}$.
\end{proof}

\begin{remark}[where the size of the lines enters]\label{rem:linesize}
Infiniteness of the lines was used only to choose a point of a line
outside a prescribed set of at most four points; so the proof works
verbatim for any plane satisfying (I1), (I2), whose point set is an
abelian group acting by collineations as in (I5), and in which every
line has at least five points.  Connes and Consani show that four points per
line suffice for the hypergroup axioms [CC11, Proposition~3.1]; on a
four-point line $L=\{a,b,c,d\}$ one computes
$a\bp b\bp c=\{0,a,b,c\}$ rather than $L\cup\{0\}$, still a symmetric
expression.
With only three points per line the rule~\eqref{eq:addition} is not
associative: on the line $\{a,b,c\}$ one finds
$(a\bp b)\bp c=c\bp c=\{0,c\}$ but $a\bp(b\bp c)=a\bp a=\{0,a\}$.
\end{remark}

\section{The counterexample}\label{sec:counterexample}

We now prove Theorem~\ref{thm:main} for the hyperfield $\Hh$ of
Proposition~\ref{prop:hyperfield} and the elements
\[
 x=1,\qquad y=t,\qquad z=t^3 .
\]
Recall that $D\supseteq A_0=\{0,1,3,20,24,29\}$ and that $-a=a$ for
every $a\in\Hh$, so the summands of the conclusion are
\[
 x^3=1,\qquad y^3=t^3,\qquad z^3=t^9,\qquad -xyz=xyz=t^4=:q .
\]

\begin{theorem}\label{thm:counterexample}
In $\Hh$, with $x=1$, $y=t$, $z=t^3$,
\[
 0\in x\bp y\bp z
 \qquad\text{but}\qquad
 0\notin x^3\bp y^3\bp z^3\bp(-xyz)\bp(-xyz)\bp(-xyz).
\]
\end{theorem}
\begin{proof}
\emph{Step 1: the hypothesis holds.}  Since $0,1,3\in D$, the three
distinct points $1,t,t^3$ lie on $L_0$.  In particular
$t^3\in L(1,t)\setminus\{1,t\}=1\bp t$, hence
\begin{equation}\label{eq:hypothesis}
 0\in t^3\bp t^3\subseteq(1\bp t)\bp t^3=x\bp y\bp z .
\end{equation}
(In fact $x\bp y\bp z=L_0\cup\{0\}$ by the collinear case in the proof
of Proposition~\ref{prop:hyperfield}, but~\eqref{eq:hypothesis} is all
we need.)

\emph{Step 2: the cubes are noncollinear, and $q$ lies on a side of
their triangle.}  Since $20,24,29\in D$,
\begin{equation}\label{eq:specialline}
 1=t^{-20+20},\quad q=t^4=t^{-20+24},\quad t^9=t^{-20+29}\ \in\ L_{-20},
\end{equation}
so $L_{-20}=L(1,t^9)$ by (I1).  On the other hand $23\notin D$:
otherwise the two distinct pairs $(23,20)$ and $(3,0)$ in $D\times D$
would both represent $3=23-20=3-0$, contradicting the uniqueness in
Lemma~\ref{lem:diff}.  Hence $t^3=t^{-20+23}\notin L_{-20}=L(1,t^9)$,
which by (I1) says precisely that $1,t^3,t^9$ are noncollinear.

\emph{Step 3: the conclusion fails.}  Put
\[
 S=x^3\bp y^3\bp z^3=(1\bp t^3)\bp t^9=T(1,t^3,t^9).
\]
By Step~2 and the triangle formula (Lemma~\ref{lem:triangle}), $S$
contains neither $0$ nor any point of the side $L(1,t^9)=L_{-20}$.
By~\eqref{eq:specialline}, $q\in L_{-20}$, so
\begin{equation}\label{eq:disjoint}
 S\cap\{0,q\}=\varnothing .
\end{equation}
Next, by~\eqref{eq:addition},
\[
 q\bp q\bp q=\{0,q\}\bp q=(0\bp q)\cup(q\bp q)=\{q\}\cup\{0,q\}=\{0,q\}.
\]
By Remark~\ref{rem:iterated} the six-term sum may be regrouped:
\[
 x^3\bp y^3\bp z^3\bp q\bp q\bp q
 =(x^3\bp y^3\bp z^3)\bp(q\bp q\bp q)=S\bp\{0,q\}.
\]
Finally, Lemma~\ref{lem:zero} together with $-a=a$ gives, for subsets
$A,B\subseteq\Hh$, that $0\in A\bp B$ if and only if
$A\cap B\neq\varnothing$.  Applied to $A=S$ and $B=\{0,q\}$,
\eqref{eq:disjoint} yields
\[
 0\notin S\bp\{0,q\}
 =x^3\bp y^3\bp z^3\bp(-xyz)\bp(-xyz)\bp(-xyz).
\]
Together with~\eqref{eq:hypothesis} this proves the theorem.
\end{proof}

Theorem~\ref{thm:counterexample} establishes Theorem~\ref{thm:main}:
$\Hh$ is countably infinite, since $P\cong\Z$, and commutative by
Proposition~\ref{prop:hyperfield}.  So the answer to the question is
\textbf{no}.

\begin{corollary}\label{cor:nonquotient}
The hyperfield $\Hh$ is not isomorphic to a quotient hyperfield $F/G$
of a field $F$ by a subgroup $G\leq F^\times$.
\end{corollary}
\begin{proof}
By Proposition~\ref{prop:quotient} every quotient hyperfield satisfies
the cubic implication, an isomorphism of hyperfields transports the
implication, and $\Hh$ violates it by Theorem~\ref{thm:counterexample}.
\end{proof}

\begin{remark}\label{rem:context}
Hyperfields that are not quotients of fields were first constructed by
Massouros~[Ma85], by a different method.  Hyperfields attached to
infinite cyclic planes built from Hall's difference sets appear in
[CC11, Example~4.7].  For quotient hyperfields $F/G$ satisfying
$a\bp a=\{0,a\}$ for all $a$ (a property our $\Hh$ shares,
by~\eqref{eq:addition}), the projective geometry attached to the
hyperaddition, whose lines are the sets $\{a,b\}\cup(a\bp b)$ for
distinct nonzero $a,b$ (for our $\Hh$ these are exactly the lines
$L(a,b)$), is Desarguesian; [KB25, Example~2.6(4) and Remark~4.3(3)]
contrasts such quotients with a hyperfield built on an infinite cyclic
group whose plane is non-Desarguesian, and thereby separates general
hyperfields from quotients of fields.
Corollary~\ref{cor:nonquotient} gives, for the specific $\Hh$
constructed here, a direct elementary verification that involves no
Desargues configuration: the cubic implication is a first-order
statement true in every $F/G$ and false in $\Hh$.
\end{remark}

\begin{remark}[what the example shows]
For distinct $x,y,z\in P$ the hypothesis $0\in x\bp y\bp z$ is the
collinearity of $x,y,z$ (by the collinear and noncollinear cases in the
proof of Proposition~\ref{prop:hyperfield}), and, since cubing is
injective on $P\cong\Z$ so that the cubes are distinct, the conclusion
holds if the cubes are collinear and otherwise holds exactly when $xyz$
avoids all three sides of the cube triangle, by the argument of
Step~3.  In a field,
or in a quotient of a field, the two conditions are linked by the
identity~\eqref{eq:cubic}.  In the cyclic plane of a perfect difference
set the two conditions are independent, and the construction exploits
this: the seed elements $0,1,3\in D$ make $x,y,z$ collinear, so the
hypothesis holds; the seed elements $20,24,29\in D$ place $xyz=t^4$ on
the side $L(1,t^9)$ of the cube triangle, so the conclusion fails; and
the uniqueness of differences built into $D$ excludes $23\in D$ (which
would duplicate the difference $3=23-20=3-0$) and so keeps the cubes
noncollinear.
\end{remark}

\section{Background and credits}\label{sec:credits}

None of the ingredients is new: the hyperfield axioms and the quotient
hyperfields $F/G$ are Krasner's [Kr83], the term canonical hypergroup
is Mittas's [Mi72], the hyperaddition ``third points of the joining
line'' is Prenowitz's [Pr43] and was combined with an incidence group
to produce hyperfields by Connes and Consani [CC11, \S3, Example~4.7],
and perfect difference sets and cyclic planes are due to Singer~[Si38]
(finite case) and Hall~[Ha47] (the group $\Z$), of whose Theorem~3.1
Lemma~\ref{lem:diff} is an explicit instance.
Proposition~\ref{prop:quotient} is an elementary observation about
quotient hyperfields.  What appears to be new is the counterexample
itself: the choice of the seed $A_0$, the resulting incidence
configuration, and the failure of the cubic implication; a literature
search did not locate the question or an earlier answer.

\medskip
\noindent\textbf{References}

\begin{description}
\item[{[CC11]}] A. Connes, C. Consani,
\emph{The hyperring of ad\`ele classes},
J. Number Theory \textbf{131} (2011), no.~2, 159--194; arXiv:1001.4260.
\item[{[Ha47]}] M. Hall, Jr.,
\emph{Cyclic projective planes},
Duke Math. J. \textbf{14} (1947), 1079--1090.
\item[{[KB25]}] P. Kowalski, P. B\l{}aszkiewicz,
\emph{The class of Krasner hyperfields is not elementary},
J. Symbolic Logic (2025), published online.
\item[{[Kr83]}] M. Krasner,
\emph{A class of hyperrings and hyperfields},
Internat. J. Math. Math. Sci. \textbf{6} (1983), no.~2, 307--311.
\item[{[Ma85]}] Ch. G. Massouros,
\emph{Methods of constructing hyperfields},
Internat. J. Math. Math. Sci. \textbf{8} (1985), no.~4, 725--728.
\item[{[Mi72]}] J. Mittas,
\emph{Hypergroupes canoniques},
Math. Balkanica \textbf{2} (1972), 165--179.
\item[{[Pr43]}] W. Prenowitz,
\emph{Projective geometries as multigroups},
Amer. J. Math. \textbf{65} (1943), 235--256.
\item[{[Si38]}] J. Singer,
\emph{A theorem in finite projective geometry and some applications to
number theory},
Trans. Amer. Math. Soc. \textbf{43} (1938), 377--385.
\end{description}

\end{document}
